\documentclass[11pt]{article}
\usepackage[margin=1in]{geometry}
\usepackage{amsmath,amssymb,amsthm,mathtools}
\usepackage{graphicx}
\usepackage{microtype}
\usepackage[hidelinks]{hyperref}

\newtheorem{theorem}{Theorem}
\newtheorem{lemma}{Lemma}
\newcommand{\R}{\mathbb R}
\newcommand{\N}{\mathsf N}
\newcommand{\ind}{\mathrel{\perp\!\!\!\perp}}

\title{A common-affine residual class and a smooth obstruction\\
for conditional entropy-power concavity}
\author{Candidate partial result for arXiv:2501.10309}
\date{19 July 2026}

\begin{document}
\maketitle

\begin{center}
\fbox{\parbox{0.91\textwidth}{\textbf{Status.}
Candidate partial result, judged likely valid. The packet gives a
non-Gaussian positive class and disproves unrestricted concavity, but it does
not characterize all admissible classes and does not address the
Ball--Nayar--Tkocz log-concave conjecture.}}
\end{center}

\section{Source question}

Fradelizi, Gavalakis, and Rapaport define, for independent random vectors
$X,Y\in\R^n$, the map
\begin{equation}\label{eq:source-map}
 F(\lambda)=
 \frac{\N_n(\sqrt\lambda X+\sqrt{1-\lambda}Y)^n}
 {\N_{n-1}(\sqrt\lambda X'+\sqrt{1-\lambda}Y')^{n-1}},
 \qquad 0\leq\lambda\leq1,
\end{equation}
where $X',Y'$ denote the first $n-1$ coordinates and
$\N_d(W)=\exp(2h(W)/d)$ \cite{FGR}. In Remark~9, equation~(18), page~7,
they ask for which classes of random variables this map is concave.

\begin{figure}[ht]
\centering
\includegraphics[width=\textwidth]{figures/open_problem_crop.png}
\caption{The source question in Remark~9 on page~7 of arXiv:2501.10309v1.}
\end{figure}

The source proves the endpoint-chord inequality
$F(\lambda)\geq\lambda F(1)+(1-\lambda)F(0)$, which is weaker than concavity
at arbitrary pairs of parameters.

\section{Idea of the proof}

The quotient in \eqref{eq:source-map} is not an opaque ratio: it is exactly
the exponential of twice a conditional entropy. If the last coordinates of
$X$ and $Y$ have the same affine regression on their first blocks and the
residuals are independent, conditioning removes the common affine term and
the problem becomes the ordinary one-dimensional entropy-power problem for
the two residuals. Costa's theorem then gives a positive result whenever one
residual is Gaussian, even though the first $n-1$ coordinates can be wholly
non-Gaussian.

The same reduction exposes an obstruction. Smooth a Rademacher variable by a
very narrow Gaussian. At unequal coefficients, four discrete modes remain
separated; at equal coefficients, the two middle modes collide. In the
small-noise limit this lowers discrete entropy by $\tfrac12\log2$, and hence
lowers entropy power by a factor tending to $1/2$. Midpoint concavity
therefore fails. The resulting random vectors are iid and have everywhere
positive smooth densities. What remains open is a broader structural
classification, especially under log-concavity.

\section{Conditional entropy reformulation}

Write $W=(W',W_n)\in\R^{n-1}\times\R$. Whenever the displayed entropies are
finite,
\begin{equation}\label{eq:conditional}
 \frac{\N_n(W)^n}{\N_{n-1}(W')^{n-1}}
 =\exp\bigl(2h(W)-2h(W')\bigr)
 =\exp\bigl(2h(W_n\mid W')\bigr).
\end{equation}
Thus $F$ is a conditional entropy power.

\section{A non-Gaussian positive class}

\begin{theorem}[Common-affine independent residuals]\label{thm:positive}
Let $X=(X',X_n)$ and $Y=(Y',Y_n)$ be independent random vectors in $\R^n$.
Assume that, for some $a\in\R^{n-1}$,
\[
 X_n=a^T X'+U,\qquad Y_n=a^T Y'+V,
\]
where $U\ind X'$ and $V\ind Y'$. Assume the finite-entropy and finite-second-
moment hypotheses needed below. Then
\begin{equation}\label{eq:residual-reduction}
 F(\lambda)=\N_1(\sqrt\lambda U+\sqrt{1-\lambda}V).
\end{equation}
In particular, if one of $U,V$ is a nondegenerate Gaussian, then $F$ is
concave on $[0,1]$.
\end{theorem}

\begin{proof}
Put $W_\lambda=\sqrt\lambda X+\sqrt{1-\lambda}Y$ and
$R_\lambda=\sqrt\lambda U+\sqrt{1-\lambda}V$. The hypotheses and the
independence of $X$ and $Y$ imply
\[
 (W_\lambda)_n=a^T(W_\lambda)'+R_\lambda,
 \qquad R_\lambda\ind (W_\lambda)'.
\]
Translation by the conditioned value and independence give
$h((W_\lambda)_n\mid(W_\lambda)')=h(R_\lambda)$. Equation
\eqref{eq:conditional} proves \eqref{eq:residual-reduction}.

Suppose, for definiteness, that $V=\tau Z$ with $Z$ standard Gaussian.
Costa's entropy-power theorem states that
$G(t)=\N_1(U+\sqrt t Z)$ is concave for $t\geq0$ \cite{Costa}. For
$0<\lambda\leq1$, homogeneity of entropy power gives
\[
 F(\lambda)
 =\lambda G\!\left(\frac{\tau^2(1-\lambda)}{\lambda}\right).
\]
The perspective $(x,y)\mapsto xG(y/x)$ of a concave function is concave on
$x>0$, $y\geq0$. Restrict it to the affine segment
$(x,y)=(\lambda,\tau^2(1-\lambda))$ and pass to the endpoints by continuity.
The case in which $U$ is Gaussian is symmetric.
\end{proof}

The first blocks $X'$ and $Y'$ in Theorem~\ref{thm:positive} are arbitrary;
in particular, the full vectors need not be Gaussian or Gaussian mixtures.

\section{Smooth iid failure of unrestricted concavity}

We use the following elementary small-noise fact.

\begin{lemma}[Finite mixtures at small noise]\label{lem:small-noise}
Let $A$ be a finite-valued real random variable with distinct support points,
and let $Z$ be an independent standard Gaussian. Then
\[
 h(A+\sigma Z)=h(\sigma Z)+H(A)+o(1)\qquad(\sigma\downarrow0),
\]
where $H$ is discrete Shannon entropy.
\end{lemma}

\begin{proof}
The mutual-information identities give
\[
 h(A+\sigma Z)-h(\sigma Z)
 =I(A;A+\sigma Z)=H(A)-H(A\mid A+\sigma Z).
\]
Let $\delta>0$ be the minimum distance between distinct support points of
$A$. Nearest-neighbor decoding has error probability at most
$\Pr\{|\sigma Z|\geq\delta/2\}$, which tends to zero. Fano's inequality
therefore gives $H(A\mid A+\sigma Z)\to0$.
\end{proof}

\begin{theorem}[Smooth iid mode-collision obstruction]\label{thm:negative}
For every $n\geq1$ there are iid random vectors $X,Y\in\R^n$ with everywhere
positive $C^\infty$ densities and finite moments of all orders for which the
map $F$ in \eqref{eq:source-map} is not concave. More explicitly, let
$S,T$ be independent Rademacher variables, let $Z_1,Z_2$ be independent
standard Gaussians, and set
\[
 U_\sigma=4S+\sigma Z_1,\qquad
 V_\sigma=4T+\sigma Z_2.
\]
For all sufficiently small $\sigma>0$, these residuals, together with
arbitrary iid independent Gaussian first blocks, give
\[
 F_\sigma(1/2)
 <\frac{F_\sigma(1/4)+F_\sigma(3/4)}2.
\]
\end{theorem}

\begin{proof}
Take $X=(G,U_\sigma)$ and $Y=(H,V_\sigma)$, where $G,H$ are iid standard
Gaussian vectors in $\R^{n-1}$, independent of all scalar variables; for
$n=1$ omit the first blocks. The vectors $X,Y$ are iid with positive smooth
densities. Theorem~\ref{thm:positive}'s reduction (which does not require a
Gaussian residual for the identity itself) yields
\[
 F_\sigma(\lambda)
 =\N_1(\sqrt\lambda U_\sigma+\sqrt{1-\lambda}V_\sigma).
\]
Because $U_\sigma,V_\sigma$ are iid,
$F_\sigma(1/4)=F_\sigma(3/4)$.

At $\lambda=1/4$, Gaussian stability gives
\[
 \sqrt\lambda U_\sigma+\sqrt{1-\lambda}V_\sigma
 \stackrel d= A+\sigma Z,
 \qquad A=2S+2\sqrt3\,T,
\]
where $Z$ is standard Gaussian. The four values of $A$ are distinct and
equiprobable, so $H(A)=2\log2$. At $\lambda=1/2$ the corresponding discrete
part is
\[
 B=2\sqrt2(S+T),
\]
whose three values have probabilities $1/4,1/2,1/4$; hence
$H(B)=\tfrac32\log2$. Lemma~\ref{lem:small-noise} now gives
\[
 \frac{F_\sigma(1/2)}{F_\sigma(1/4)}
 =\exp\!\left(2[h(B+\sigma Z)-h(A+\sigma Z)]\right)
 \longrightarrow
 \exp(-\log2)=\frac12.
\]
Thus $F_\sigma(1/2)<F_\sigma(1/4)=F_\sigma(3/4)$ for all sufficiently small
$\sigma$, contradicting midpoint concavity.
\end{proof}

\section{Computational sanity check}

The script \texttt{code/explore\_gaussian\_mixtures.py} numerically integrates
the exact four-component Gaussian-mixture density. With component means
$\pm4$ and component standard deviation $\sigma=0.1$, an 80,001-point grid
gave
\[
 (F(1/4),F(1/2),F(3/4))
 \approx(2.7327149513,1.3663574756,2.7327149513).
\]
This is only a check. The proof of Theorem~\ref{thm:negative} is the
small-noise argument above and does not rely on numerical quadrature or on
proving that $0.1$ lies below an explicit threshold.

\section{Verification and limitations}

An adversarial step check is supplied in \texttt{verification.md}. The core
proof is self-contained apart from Costa's theorem in the positive direction.
The negative theorem uses only the entropy chain rule, elementary decoding,
and Fano's inequality.

The result is partial. It neither classifies all pairs $(X,Y)$ nor resolves
the log-concave two-weight entropy conjecture cited by the source. The
negative examples are smooth and symmetric but non-log-concave for small
$\sigma$.

\paragraph{Bounded novelty check.}
Searches made on 19 July 2026 used the exact source phrase and combinations of
entropy power concavity, Rademacher Gaussian mixture, iid, and Costa. The
search also checked the generalized-EPI counterexample paper of
Madiman--Nayar--Tkocz and the centered Gaussian scale-mixture paper of
Eskenazis--Gavalakis \cite{MNT,EG}. Those works concern related entropy or
multiweight phenomena. No exact match was found for the conditional quotient
reduction or the two-weight smooth mode-collision proof above. This was a
bounded search, not an exhaustive novelty certification.

\paragraph{Human-review recommendation.}
Send to a specialist reviewer. The reviewer should check the preferred
regularity formulation of Costa's theorem, the independence bookkeeping in
Theorem~\ref{thm:positive}, and whether a later version of the source already
records the same common-affine residual corollary.

\small
\begin{thebibliography}{9}
\bibitem{FGR}
M.~Fradelizi, L.~Gavalakis, and M.~Rapaport,
\emph{Entropic versions of Bergstr\"om's and Bonnesen's inequalities},
arXiv:2501.10309v1, 2025.

\bibitem{Costa}
M.~H.~M. Costa,
\emph{A new entropy power inequality},
IEEE Trans. Inform. Theory 31 (1985), no.~6, 751--760.

\bibitem{MNT}
M.~Madiman, P.~Nayar, and T.~Tkocz,
\emph{Two remarks on generalized entropy power inequalities},
arXiv:1904.02314, 2019.

\bibitem{EG}
A.~Eskenazis and L.~Gavalakis,
\emph{On the entropy and information of Gaussian mixtures},
Mathematika 70 (2024), e12246; arXiv:2308.15997.
\end{thebibliography}

\end{document}
